\chapter{总结与展望}

\section{线性注意力的核心贡献}

线性注意力机制通过将标准 softmax 注意力重写为可分解的核函数形式，成功地将注意力计算的时间复杂度从 $\mathcal{O}(N^2 d)$ 降低到 $\mathcal{O}(Nd^2)$，空间复杂度从 $\mathcal{O}(N^2)$ 降低到 $\mathcal{O}(Nd)$。这使得 Transformer 类模型能够处理更长的序列。

线性注意力的核心思想可以概括为以下三个步骤：

\begin{enumerate}
    \item 将 softmax 核函数替换为可分解的核函数 $\kappa(\mathbf{q}, \mathbf{k}) = \phi(\mathbf{q})^\top \phi(\mathbf{k})$
    \item 利用矩阵乘法的结合律，将计算顺序从 $(\mathbf{Q}\mathbf{K}^\top)\mathbf{V}$ 改为 $\mathbf{Q}(\mathbf{K}^\top\mathbf{V})$
    \item 通过累积聚合量 $\mathbf{S} = \sum \phi(\mathbf{k}_j)\mathbf{v}_j^\top$ 和 $\mathbf{z} = \sum \phi(\mathbf{k}_j)$ 实现线性复杂度
\end{enumerate}

\section{主要方法回顾}

\begin{table}[H]
    \centering
    \begin{tabular}{l|l|c|c}
        \toprule
        方法 & 特征映射 & 时间复杂度 & 关键特性 \\
        \midrule
        Linear Transformer & $\text{elu}(\mathbf{x}) + 1$ & $\mathcal{O}(Nd^2)$ & 开创性工作，简单高效 \\
        Performer & 随机傅里叶特征 & $\mathcal{O}(Nmd)$ & 理论保证，接近 softmax \\
        GLA & elu+1 + 门控 & $\mathcal{O}(Nd^2)$ & 门控增强信息控制 \\
        RWKV & 时间衰减 & $\mathcal{O}(Nd)$ & RNN 推理效率 \\
        RetNet & 多尺度衰减 & $\mathcal{O}(Nd^2)$ & 训练推理双模式 \\
        \bottomrule
    \end{tabular}
    \caption{主要线性注意力方法总结}
\end{table}

\section{与 Mamba/SSM 的关系}

近年来，Mamba 等状态空间模型（SSM）也在长序列建模领域取得了显著进展。线性注意力与 SSM 有密切的联系：

\begin{itemize}
    \item 两者都具有 $\mathcal{O}(N)$ 或 $\mathcal{O}(Nd)$ 的递推形式
    \item 线性注意力的状态矩阵 $\mathbf{S}_t$ 类似于 SSM 的隐藏状态
    \item 最新的 Mamba-2 工作表明，SSM 可以被视为一种特殊的线性注意力
    \item 两者的界限正在逐渐模糊，出现了统一的视角
\end{itemize}

\section{未来方向}

线性注意力的研究仍有以下值得探索的方向：

\begin{enumerate}
    \item \textbf{更好的特征映射}：设计更接近 softmax 且计算高效的特征映射，减少表达能力差距
    \item \textbf{理论分析}：深入理解线性注意力的表达能力边界和最优特征映射
    \item \textbf{混合架构}：探索标准注意力与线性注意力的最优组合方式
    \item \textbf{硬件优化}：针对线性注意力的计算模式设计专用的硬件加速方案
    \item \textbf{多模态应用}：将线性注意力扩展到视觉、音频、多模态等领域
    \item \textbf{统一框架}：建立线性注意力、SSM 和其他高效注意力方法的统一理论框架
\end{enumerate}

\section{结语}

线性注意力是大模型高效架构设计的重要方向之一。随着理论认识的深入和新方法的出现，线性注意力有望在保持高效的同时，进一步缩小与标准注意力的性能差距，成为处理超长序列的标配方案。
